\section{Introduction}
% Overview figure (first page). Snapshots: scripts/build_journal_snapshots.py.
% Type floor: \scriptsize (7 pt in IEEEtran journal 10 pt); nothing smaller. Change maps live in Fig. qualitative.
\begin{figure*}[t]
\centering
\begin{tikzpicture}[
  font=\scriptsize, >={Stealth[length=1.6mm]},
  panel/.style={draw=snfgrey!60,rounded corners=2pt,fill=snfpaper,inner sep=3pt},
  head/.style={font=\footnotesize\bfseries,text=snfink,anchor=south west,inner sep=1pt},
  img/.style={inner sep=0pt,outer sep=0pt},
  stage/.style={draw=snfink!70,fill=white,rounded corners=1.5pt,minimum height=11mm,
    text width=3.15cm,align=center,font=\footnotesize,inner sep=2pt},
  flow/.style={->,line width=.7pt,snfink!80}]

% ---- (a) inputs
\node[head] at (0,4.6) {(a) Inputs};
\node[panel,anchor=north west,text width=3.4cm,align=left] (txt) at (0,4.45)
  {\textbf{T2V}\enspace\textit{``\dots recorded by a fixed tripod camera. Rain falls continuously\dots''}};
\node[img,anchor=north west] (src) at ($(txt.south west)+(0,-0.15)$)
  {\includegraphics[width=3.6cm]{figures/snapshots/role_piq0140.png}};
\node[anchor=north west,align=left,inner sep=1pt] at ($(src.south west)+(0,-0.05)$)
  {\textbf{I2V} source with role map:\\
   \swatch{roleSupport}\,support\quad\swatch{roleMaterial}\,material\\
   \swatch{roleIgnore}\,uncertain (enters no measure)};

% ---- (b) long rollout of one published output
\node[head] at (4.05,4.6) {(b) One 240\,s rollout, fixed-camera prompt};
\foreach \t/\x in {000/0,060/1,120/2,240/3} {
  \node[img,anchor=north west] (r\x) at ({4.1+1.9*\x},4.4)
    {\includegraphics[width=1.84cm]{figures/snapshots/t2v240/p0/CausVid_t\t.png}};
}
\draw[->,line width=.6pt,snfink] ($(r0.south west)+(0,-0.12)$) -- ($(r3.south east)+(0.12,-0.12)$)
  node[right] {$t$};
\foreach \lab/\x in {0/0,60/1,120/2,240/3}
  \node[below] at ($(r\x.south)+(0,-0.12)$) {\gclock\,\lab\,s};
\node[text=snfink!75,anchor=north west,align=left,text width=7.4cm,inner sep=1pt] at ($(r0.south west)+(0,-0.68)$)
  {The layout departs from the first frame despite the fixed-camera prompt. Whole-frame change mixes
   intended and unintended effects; Fig.~\ref{fig:qualitative} adds change maps for all eight checkpoints.};

% ---- (c) measures with status
\node[head] at (12.0,4.6) {(c) Measures and status};
\node[panel,anchor=north west,text width=5.75cm,align=left] (meas) at (12.0,4.45) {%
  \setlength{\tabcolsep}{1.5pt}\renewcommand{\arraystretch}{1.12}%
  \begin{tabular}{@{}c@{\hspace{3pt}}l@{\hspace{5pt}}l@{\hspace{4pt}}c@{}}
  \gsupport   & Support fidelity     & fBD\down\,NBF\down; DLR ctx. & \stfrozen\\
  \gmotion    & Material motion      & MCFF, FP; DAR ctx.           & \stfrozen\\
  \greplay    & Replay\,/\,progression & R4                         & \stcand\\
  \gdirection & Direction            & R3                           & \stcand\\
  \gsupport   & Support integrity    & R1                           & \stdiag\\
  \gmaterial  & Material presence    & R2                           & \stopen\\
  \gplaus     & Motion plausibility  & R5                           & \stopen\\
  \gintegrity & Scene integrity      & ---                          & \stopen\\
  \ghuman     & Human validation     & ---                          & \stopen
  \end{tabular}};
\node[anchor=north west,align=left,inner sep=1pt] at ($(meas.south west)+(0,-0.08)$)
  {\stfrozen\,frozen definition\quad\stcand\,development candidate\\
   \stdiag\,diagnostic only\quad\stopen\,planned\quad ctx.\,=\,context only};

% ---- stages
\node[stage] (s1) at (1.645,-0.55) {\cmark\ Specifications\\[-1pt]{\scriptsize 96 T2V + 48 I2V\\nine regimes}};
\node[stage,right=3mm of s1] (s2) {\stcand\ Role maps\\[-1pt]{\scriptsize automatic (I2V)\\2 annotators + adjudicator}};
\node[stage,right=3mm of s2] (s3) {\stcand\ Measures\\[-1pt]{\scriptsize versioned identities\\explicit abstentions}};
\node[stage,right=3mm of s3] (s4) {\stcand\ Development\\[-1pt]{\scriptsize matched controls\\real footage}};
\node[stage,right=3mm of s4,fill=snfpaper] (s5) {\glock\ Confirmation\\[-1pt]{\scriptsize held-out scenes\\and model families}};
\foreach \a/\b in {s1/s2,s2/s3,s3/s4,s4/s5} \draw[flow] (\a) -- (\b);
\node[text=snfink!75,anchor=north] at ($(s3.south)+(0,-0.06)$)
  {\ghuman\ blinded human ratings per construct feed development and confirmation};
\end{tikzpicture}
\caption{\textbf{SNF-Bench++ at a glance.} (a)~Each scene specifies a fixed viewpoint and names a natural process; I2V
sources carry role maps; illustrated on a public fixed-camera scene, smoke rising from burning paper, with roles derived from that benchmark's motion masks (overlay pixels such as falling snow, also excluded, are absent here). (b)~Frames of a historical output of a published checkpoint~\cite{yin2025causvid} under the common inference configuration for a
fixed-camera prompt (first system alphabetically, first 240\,s prompt; Fig.~\ref{fig:qualitative}).
(c)~Measures by construct and validation status; no summed score.
Bottom: the protocol, from specification to held-out confirmation. Image in~(a): Physics-IQ benchmark~\cite{motamed2025physicsiq}, scene 0140, first frame, \copyright~2024 DeepMind Technologies Limited, \href{https://creativecommons.org/licenses/by/4.0/}{CC~BY~4.0}; resized, with roles overlaid. An interactive version accompanies the
supplementary material (\texttt{interactive/measurement\_map.html}).}
\label{fig:pp_protocol}
\end{figure*}

\IEEEPARstart{I}{n} a fixed-camera video of a river, the water should move while the banks stay in place. Long generated videos can fail in either direction: the landscape may drift while the river moves, or the image may stay stable because the river has frozen. A useful evaluation must show both failures separately.

Existing benchmarks already measure important aspects of video quality, motion and consistency~\cite{huang2023vbench,huang2024vbench++,Zheng2025VBench20AV,liu2024evalcrafter,duan2025worldscore}. Our question is whether a fixed-camera evaluation adds useful information on the same videos, not whether those benchmarks are unable to see the problem. The distinction matters at long horizons, where small changes can accumulate over minutes~\cite{huang2025selfforcing,liu2025rolling,yang2025longlive,helios,li2025stable}.

SNF-Bench++ evaluates text-to-video and image-to-video generation for fixed-camera scenes with continuous natural motion. It separates the \emph{background}, which should remain stable, from the \emph{material}, such as water or smoke, that the prompt asks to move. We report background displacement and motion alongside early and late material motion. We also study checks for motion location, direction, repetition and plausibility. These checks are candidates until their controls, real-footage comparisons and human studies are complete.

The original SNF-Bench preprint~\cite{minar2026snfbench} and the source-fixed method-paper protocol are prior work. Appendix~\ref{app:prior} will distinguish their contributions from this journal study. Human and animal subjects are possible extension tracks; they are not evidence for the present nature-flow task.

This working manuscript sets out:
\begin{itemize}
\item \textbf{A scene suite} for both generation tracks, with planned 60, 120 and 240 second horizons and documented region maps (Section~\ref{sec:benchmark}).
\item \textbf{Separate measurements} of background change, material motion and temporal behavior, with explicit limits and abstentions (Section~\ref{sec:measurement}).
\item \textbf{A validation plan} using targeted controls, real footage, human judgments and comparisons on shared inputs (Section~\ref{sec:validation}).
\item \textbf{An evaluation plan} for released long-video generators (Section~\ref{sec:results}). \pending{independent panel, human validation and headline findings}
\end{itemize}
